\documentclass[10pt,a4paper]{article}

\usepackage{geometry}
\geometry{
  left=15mm,
  right=15mm,
  top=15mm,
  bottom=15mm
}
\usepackage{amsmath}
\usepackage{amssymb}
\usepackage{enumitem}

\setlist[enumerate,1]{label=(\arabic*),leftmargin=2.2em}
\setlist[enumerate,2]{label=(\alph*),leftmargin=2.2em}

\begin{document}

\begin{center}
  {\Large \textbf{The University of Tokyo 2022}}\\[2mm]
  {\large \textbf{Problems 1--3}}
\end{center}

\section*{Problem 1}

Consider the following multiple conditions on $x,y,z\in\mathbb{R}$:
\[
  \begin{cases}
    0<z-xy<1,\\
    0<z-(x+y)^2<-xy.
  \end{cases}
\]
Let $\Omega$ be the set of points $(x,y)$ for which at least one $z$ exists satisfying the above conditions. Note that the set $\Omega$ can be seen in the three-dimensional Cartesian coordinate system as the orthogonal projection of points $(x,y,z)$ satisfying the above conditions onto the $xy$-plane. Answer the following questions.

\begin{enumerate}
  \item Find the inequalities on $x$ and $y$ representing $\Omega$.

  \item Draw a figure of $\Omega$ in the $xy$-plane. If the boundary of $\Omega$ intersects with the $x$-axis or the $y$-axis, write down the coordinates at each intersection.

  \item The curved segments of the boundary of $\Omega$ correspond to the linear transformation of arcs of the unit circle with a matrix $X$. Find one such $X$. Note that the point $(1,0)$ on the unit circle must be transformed to a point where the curvature is maximized in the curved segments.

  \item Calculate the determinant of $X$ found in (3).

  \item Calculate the area of the set $\Omega$. Note that the absolute value of the determinant of a matrix is the area scale factor of the transformation with that matrix.
\end{enumerate}

\section*{Problem 2}

Consider the following integral $I_n(\alpha)$ for $\alpha\geq 1$ and $n>0$:
\[
  I_n(\alpha)=\int_{1/n}^{n}\frac{f(\alpha x)-f(x)}{x}\,dx.
\]
Assume that a real-valued function $f(x)$ is continuous and differentiable on $x\geq 0$, its derivative is continuous, and
\[
  \lim_{x\to\infty}f(x)=0.
\]
Answer the following questions.

\begin{enumerate}
  \item Define
        \[
          J_n(\alpha)=\frac{dI_n(\alpha)}{d\alpha}.
        \]
        Show that
        \[
          J_n(\alpha)=\frac{1}{\alpha}
          \left\{f(\alpha n)-f\left(\frac{\alpha}{n}\right)\right\}.
        \]
        You can use the fact that the integration and the differentiation commute in this context.

  \item Define
        \[
          I(\alpha)=\lim_{n\to\infty}I_n(\alpha).
        \]
        Show that $\lim_{n\to\infty}J_n(\beta)$ exists for any $\beta\in[1,\alpha]$ and it uniformly converges on $[1,\alpha]$, and show that
        \[
          I(\alpha)=\int_1^\alpha
          \left(\lim_{n\to\infty}J_n(\beta)\right)\,d\beta.
        \]

  \item Obtain $I(\alpha)$.

  \item Calculate the following integral. Note that $p>q>0$.
        \[
          \int_0^\infty
          \frac{e^{-px}\cos(px)-e^{-qx}\cos(qx)}{x}\,dx.
        \]
\end{enumerate}

\section*{Problem 3}

Consider a region $R$ defined by $0<x<1$ and $0<y<1$ in the $xy$-plane. We randomly select a point on $R$ and refer to the selected point as A. We assume that A is uniformly distributed on $R$. Let AB be a perpendicular line from A to the $y$-axis and AC be a perpendicular line from A to the $x$-axis as shown in the figure. We call rectangle OCAB as the rectangle of A, where O denotes the origin. Let $S$ be a random variable representing the area of the rectangle of A. Answer the following questions.

\begin{enumerate}
  \item Calculate the expectation value of $S$.

  \item Calculate the probability that $S\leq r$ holds, where $0<r<1$.

  \item Calculate the probability density function of $S$.
\end{enumerate}

Again consider the region $R$. Let $n$ be a positive integer. We select $n$ points on $R$ and refer to the selected points as $A_1,A_2,\ldots,A_n$. We assume that each of the points is uniformly distributed on $R$, and $A_i$ and $A_j$ for $i\neq j$ are selected independently. Answer the following question.

\begin{enumerate}[resume]
  \item Let $S_i$ be a random variable representing the area of the rectangle of $A_i$. Let $Z$ be a random variable which is the minimum of $S_1,S_2,\ldots,S_n$. Calculate the probability density function of $Z$.
\end{enumerate}

\begin{center}
\setlength{\unitlength}{0.8mm}
\begin{picture}(100,70)
  \put(15,10){\vector(1,0){70}}
  \put(15,10){\vector(0,1){55}}
  \put(87,7){$x$}
  \put(18,66){$y$}
  \put(15,10){\line(1,0){55}}
  \put(15,10){\line(0,1){45}}
  \put(70,10){\line(0,1){45}}
  \put(15,55){\line(1,0){55}}
  \multiput(18,12)(5,0){10}{\line(1,1){12}}
  \put(67,57){A}
  \put(10,52){B}
  \put(67,5){C}
  \put(10,5){O}
  \put(42,31){$S$}
  \dashbox{1}(55,45){}
  \put(67,57){\circle*{1.5}}
  \put(67,62){$(1,1)$}
  \put(4,55){$(0,1)$}
  \put(63,1){$(1,0)$}
\end{picture}
\end{center}

\newpage

\begin{center}
  {\Large \textbf{Solutions 1--3}}
\end{center}

\section*{Solution 1}

\begin{enumerate}
  \item For a fixed $(x,y)$, the first condition says
        \[
          xy<z<xy+1,
        \]
        and the second condition says
        \[
          (x+y)^2<z<(x+y)^2-xy.
        \]
        Thus such a $z$ exists if and only if the two open intervals
        \[
          (xy,xy+1)
          \quad\text{and}\quad
          ((x+y)^2,(x+y)^2-xy)
        \]
        overlap.

        The second interval is non-empty precisely when
        \[
          xy<0.
        \]
        Also, the two intervals overlap precisely when
        \[
          (x+y)^2<xy+1.
        \]
        This is equivalent to
        \[
          x^2+xy+y^2<1.
        \]
        Therefore
        \[
          \boxed{
          \Omega=\{(x,y)\in\mathbb{R}^2\mid xy<0,\ x^2+xy+y^2<1\}.}
        \]

  \item The set $\Omega$ consists of the two parts of the ellipse
        \[
          x^2+xy+y^2<1
        \]
        lying in the second and fourth quadrants. The boundary consists of the curved ellipse segments together with the portions of the coordinate axes in the closed ellipse.

        The boundary intersects the $x$-axis on
        \[
          \boxed{\{(x,0)\mid -1\leq x\leq 1\}},
        \]
        and it intersects the $y$-axis on
        \[
          \boxed{\{(0,y)\mid -1\leq y\leq 1\}}.
        \]
        In particular, the ellipse itself meets the axes at
        \[
          \boxed{(\pm1,0),\ (0,\pm1).}
        \]

        \begin{center}
        \setlength{\unitlength}{0.8mm}
        \begin{picture}(110,80)
          \put(55,40){\vector(1,0){45}}
          \put(55,40){\vector(-1,0){45}}
          \put(55,40){\vector(0,1){32}}
          \put(55,40){\vector(0,-1){32}}
          \put(101,37){$x$}
          \put(58,72){$y$}
          \qbezier(15,70)(35,52)(55,40)
          \qbezier(55,40)(75,28)(95,10)
          \qbezier(15,70)(35,88)(55,40)
          \qbezier(55,40)(75,-8)(95,10)
          \put(30,57){$\Omega$}
          \put(76,20){$\Omega$}
          \put(80,43){$(1,0)$}
          \put(17,36){$(-1,0)$}
          \put(58,64){$(0,1)$}
          \put(58,13){$(0,-1)$}
        \end{picture}
        \end{center}

  \item The ellipse is
        \[
          \begin{bmatrix}x&y\end{bmatrix}
          \begin{bmatrix}
            1 & 1/2\\
            1/2 & 1
          \end{bmatrix}
          \begin{bmatrix}x\\y\end{bmatrix}
          =1.
        \]
        Its major-axis endpoints on the curved boundary are
        \[
          (1,-1)
          \quad\text{and}\quad
          (-1,1),
        \]
        where the curvature is maximized. One suitable matrix that maps the unit circle to this ellipse and maps $(1,0)$ to $(1,-1)$ is
        \[
          \boxed{
          X=
          \begin{bmatrix}
            1 & 1/\sqrt{3}\\
            -1 & 1/\sqrt{3}
          \end{bmatrix}.}
        \]

  \item For the matrix above,
        \[
          \det X
          =
          \frac{1}{\sqrt3}+\frac{1}{\sqrt3}
          =
          \boxed{\frac{2}{\sqrt3}.}
        \]

  \item Under the transformation $X$, the condition $xy<0$ becomes
        \[
          |v|<\sqrt3|u|
        \]
        in the unit disk with coordinates $(u,v)$. In polar coordinates this corresponds to total angle
        \[
          \frac{4\pi}{3}.
        \]
        Hence the corresponding area in the unit disk is
        \[
          \frac12\cdot\frac{4\pi}{3}
          =
          \frac{2\pi}{3}.
        \]
        Multiplying by the area scale factor $|\det X|=2/\sqrt3$, we obtain
        \[
          \boxed{
          \operatorname{Area}(\Omega)=
          \frac{4\pi}{3\sqrt3}.}
        \]
\end{enumerate}

\section*{Solution 2}

\begin{enumerate}
  \item Differentiating under the integral sign gives
        \[
          J_n(\alpha)
          =
          \int_{1/n}^{n} f'(\alpha x)\,dx.
        \]
        With $u=\alpha x$, $du=\alpha\,dx$, so
        \[
          J_n(\alpha)
          =
          \frac1\alpha
          \int_{\alpha/n}^{\alpha n}f'(u)\,du
          =
          \boxed{
          \frac1\alpha
          \left\{f(\alpha n)-f\left(\frac{\alpha}{n}\right)\right\}.}
        \]

  \item For $\beta\in[1,\alpha]$,
        \[
          J_n(\beta)
          =
          \frac1\beta
          \left\{f(\beta n)-f\left(\frac{\beta}{n}\right)\right\}.
        \]
        Since $\beta n\geq n$ and $\lim_{x\to\infty}f(x)=0$,
        \[
          f(\beta n)\to 0
        \]
        uniformly on $[1,\alpha]$. Also, since $0\leq \beta/n\leq \alpha/n$ and $f$ is continuous at $0$,
        \[
          f(\beta/n)\to f(0)
        \]
        uniformly on $[1,\alpha]$. Hence
        \[
          J_n(\beta)\to -\frac{f(0)}{\beta}
        \]
        uniformly on $[1,\alpha]$.

        Since $I_n(1)=0$, we have
        \[
          I_n(\alpha)=\int_1^\alpha J_n(\beta)\,d\beta.
        \]
        The uniform convergence allows the limit to pass through the integral, so
        \[
          \boxed{
          I(\alpha)=\int_1^\alpha
          \left(\lim_{n\to\infty}J_n(\beta)\right)d\beta.}
        \]

  \item From part (2),
        \[
          I(\alpha)
          =
          \int_1^\alpha -\frac{f(0)}{\beta}\,d\beta
          =
          \boxed{-f(0)\log\alpha.}
        \]

  \item Put
        \[
          f(x)=e^{-x}\cos x,
          \qquad
          \alpha=\frac{p}{q}.
        \]
        Then $f(0)=1$ and, using $u=qx$,
        \[
          \int_0^\infty
          \frac{e^{-px}\cos(px)-e^{-qx}\cos(qx)}{x}\,dx
          =
          \int_0^\infty
          \frac{f(\alpha u)-f(u)}{u}\,du.
        \]
        Therefore, by part (3),
        \[
          \boxed{
          \int_0^\infty
          \frac{e^{-px}\cos(px)-e^{-qx}\cos(qx)}{x}\,dx
          =
          -\log\frac{p}{q}
          =
          \log\frac{q}{p}.}
        \]
\end{enumerate}

\section*{Solution 3}

\begin{enumerate}
  \item If $A=(X,Y)$, then $X$ and $Y$ are independent uniform random variables on $(0,1)$, and
        \[
          S=XY.
        \]
        Therefore
        \[
          E[S]=E[X]E[Y]=\frac12\cdot\frac12.
        \]
        Hence
        \[
          \boxed{E[S]=\frac14.}
        \]

  \item For $0<r<1$,
        \[
          P(S\leq r)
          =
          P(XY\leq r).
        \]
        Since $0<X<1$ and $0<Y<1$,
        \[
          P(XY\leq r)
          =
          \int_0^r1\,dx+\int_r^1\frac{r}{x}\,dx.
        \]
        Thus
        \[
          P(S\leq r)
          =
          r+r\int_r^1\frac{1}{x}\,dx
          =
          \boxed{r(1-\log r).}
        \]

  \item The distribution function is
        \[
          F_S(r)=r(1-\log r)
          \qquad (0<r<1).
        \]
        Hence the density is
        \[
          \boxed{
          f_S(r)=-\log r
          \qquad (0<r<1).}
        \]
        Also $f_S(r)=0$ outside $(0,1)$.

  \item Since $Z=\min(S_1,\ldots,S_n)$ and the $S_i$ are independent,
        \[
          P(Z>z)=P(S_1>z,\ldots,S_n>z)
          =
          \{1-F_S(z)\}^n.
        \]
        Thus
        \[
          F_Z(z)=1-\{1-z(1-\log z)\}^n
          =
          1-\{1-z+z\log z\}^n.
        \]
        Differentiating gives
        \[
          \boxed{
          f_Z(z)=
          n(-\log z)\{1-z+z\log z\}^{n-1}
          \qquad (0<z<1).}
        \]
        Also $f_Z(z)=0$ outside $(0,1)$.
\end{enumerate}

\end{document}
